\documentclass[a4paper,11pt]{article}

\usepackage{amsmath,amsfonts}
\usepackage{hyperref}
\usepackage[margin=2cm]{geometry}

\newcommand\lean{\texttt}
\newcommand\LEAN{\textsc{lean}}
\newcommand\MW{Minkowski-Weyl}

\begin{document}

\section{H-representation versus V-representations}
\label{sec:h-versus-v}



Various types of polyhedral objects (polytopes, polyhedra, polyhedral
cones)
can be described in two alternative ways:
A \emph{V-representation} represents them as the hull
% (convex hull or positive hull, resp
of some generating elements (points, rays), while
an \emph{H-representation} represents them as intersection of half-spaces.

The Minkowski--Weyl family of theorems says that these two types of
representations define more or less the same classes of objects
(sometimes under extra conditions).


When the workshop started, the \MW\ theorem for polyhedral cones was already in
place.

For cones, there is duality operation, which
exchanges
V-representations with H-representations.
Therefore, the \MW\ theorem can be
equivalently expressed as saying that
the dual of a V-cone is a
V-cone, 
and the dual of an H-cone is an H-cone.


\paragraph{Generality.}
The general motto is to keep results as general as possible,
if the effort is not too big.
Therefore the results are formulated in such a way that they 
work in infinite-dimensional spaces:
In the H-representation, the given half-spaces are intersected with an
affine
subspace (of arbitrary dimension), and the V-representation can use
an 
additional linear
subspace (of arbitrary dimension)
to span the objects.

\subsection{Progress}
\label{sec:progress}

\paragraph{Homogenization.}
The standard approach
for ``lifting'' the \MW\ Theorem from cones to polytopes and polyhedra
% to transferring these results to polyhedra
is through \emph{homogenization}, embedding the
affine space in which the objects lie into a (linear) vector space of
one dimension higher.
A framework for homogenization exists already in Mathlib.


The straightforward approach of homogenizing a polyhedron to a
polyhedral cone fails, because the cone spanned by an affine ray is not
closed (in the topological sense) and therefore cannot be directly used.

One solution is to explicitly add the recession cone to the
homogenization
(either as a disjoint set union, by Minkowski summation).

Another solution, starting from an H-polyhedron, simply lifts the
affine half-spaces to linear half-spaces (extending them through the
origin) and
adds the linear constraint that the extra coordinate
(``\lean{weight}'') is nonnegation.

Both approaches lead to the same polyhedral cone.
The second approach
never explicitly mentions the recession cone.


\subsection{Terminology}
\label{sec:terminology}

All objects exist in H and V variations, depending on how they are
given.

these objects are

For algorithmic questions and effective operations on polytopes and
polyhedra, it is important, which data are available.
On the other hand,
the \MW\ theorems say that there is no mathematical distiction.

After some discussion, it was decided
that \lean{Polytope}/\lean{IsPolytope} should be defined as V-polytope,
and \lean{Polyhedron}/\lean{IsPolyhedron} should be defined as
H-polyhedron,

This is
in accordance with the standard usage that a polytope, as a bounded object,
arises naturally as the convex hull of finitely many points, % vertices,
whereas the intersection of half-spaces may be unbounded and give rise to
 a polyhedron.

In the realm of cones,
\lean{IsPolyhedral} is already defined as a V-cone.

H-polyhedron is defined as the
intersection of finitely many halfspaces with an affine subspace
(which can be inifinite-dimensional).

$P
\mathop{+_\mathrm{a}}
C \mathop{+_\mathrm{v}}
S
$, where P is the convex hull of finitely many points,
$C$ is a finitely generated cone
and $S$ is a linear subspace (\lean{Submodule} in mathlib
terminology),
which can be inifinite-dimensional.



 H and V variations, depending on how they are


\subsubsection{Recession cones}
\label{sec:recession}

Part of the group worked independently on recession cones of
polyhedra, which, as it turns out, have some properties that
do not hold for general convex sets. For example,
\lean{x.recess
  $\mathop{+%_v
  }$ y.recess =
(x
$\mathop{+%_v
}$ y).recess} does not hold
for general convex sets,
where $+$ denotes Minkowski addition.
(The inclusion in the
$\subseteq$ direction holds in general.)

A difficulty is that polytopes and polyhedra live in affine spaces.
In affine space, one can add a point and a vector, but one cannot add
two points.

\subsection{Wishlist}
\label{sec:wishlist}

\begin{enumerate}
\item 
 Every bounded polyhedron is a polyope (that is, without cone and subspace).
 This should be an easy consequence.

 One challenge is to define boundedness without introducing a norm.
\item faces of polyhedra (with connection to the group that is working
  on edge graphs of polytopes)
\item polarity
   (with connection to the group that is working
  on operations on polytopes)
\item Uniqueness of the V-representation: There is a minimal set of
  vertices and extreme rays that span a given polyhedron, which is
  unique up to adding element of the lineality space.
\end{enumerate}




\subsection{Difficulties}
\label{sec:difficult}

\subsubsection{Mathematical Difficulties}
\label{sec:diff-math}

\subsubsection{Difficulties with Mathlib and \LEAN}
\label{sec:diff-lean}

Mathlib uses the term \lean{PointedCone} to denote a cone that must
contain $0$ (which is not a necessary condition for a \lean{Cone} in general).
On the other hand,
in polyhedral theory a \emph{pointed cone} is a convex cone that has a vertex
at the origin, or in other words, that has
only a trivial lineality space. Such a cone is called a 
\lean{SalientCone} in Mathlib.
We heard that efforts are under way to change the names in Mathlib.


\lean{hom.weight x} and the shorter notation
\lean{x.weight}, the latter referring to the \emph{canonical
  homogenization}.
Since there was no reason to choose any particular 
homogenization, the proof uses the canonical
homogenization

The two expressions should be definitionally equal.


extra proof steps to provide glue between the two notions.



transition

\lean{ConvexSet.convexhull}
\lean{Convexity.convexhull}

(one of them produces a \lean{Set}
the other a \lean{Pointedcone})


painful experience for the newcomers to \LEAN.

\end{document}
